\section{Introduction}
Proving a committed string belonging (or not) to a regular expression
specification has wide applications. For example, it is shown
in Zombie \cite{Zombie23} and ZkMiddleBox \cite{ZkMiddleBox}
how to prove in zero knowledge an encrypted network packet
does not contain addresses in a blocked list. In \cite{Zkreg}, 
all Linux executables are proved to be
malware free regarding the ClamAV hex-signatures.  
In \cite{Reef}, committed DNA sequences are shown to contain
chromosomes. In \cite{zkemail}, zero knowledge proof is
applied to redacted emails, twitters, and credit records.


Various techniques have different aims, e.g., striving for
expressiveness (providing negation in Zombie \cite{Zombie23}), 
efficiency for large collections of fixed patterns  \cite{Zkreg},
skipping widecard matches for handling very long input strings \cite{Reef},
and general NFA \cite{Luo23}. 

This paper considers two research problems: (1) can we
scale variety of zkreg proof techniques for handling a large
{\em collection} of regex-patterns and prove that a hidden string
{\em is not a member} of any of these patterns? 
Note that this is a difficult problem,
except for the case of fixed patterns using AC-DFA \cite{Zkreg}, 
to intersect or negate regex can easily lead to exponential cost; 
(2) does there exist more efficient weaker regex subset
which leads to more efficient techniques?
We tackle the above research problems by providing the
following contributions.
We find our technique very useful, e.g., in policy compliant
secure messaging \cite{alwen2025policy}, where content
moderation is needed in certain context, such as preventing child sexual abuse.

{\bf Empirical analysis of regex in practice:} 
The analysis forms the basis for efficient ZKP techniques we 
introduce in this paper. 
We perform an empirical analysis of the clamav regex 
and a comprehensive data that that comprises of
all executables in a Linux CentOS, malware samples,
both benign and malicious emails and SQL database logs (TODO).
We observe that for all benign binary executables in Linux, 
in average they have no more than $4$ ``near-misses", i.e., containing
any of the fixed string patterns extracted from ClamAV regex signatures.
Similar observations apply to benign email and SQL database logs,
which implies the effectiveness for a disjunctive ZKP scheme
which only needs the prover to pay for cost of discharging
those ``near-miss" regex signatures.
We also notice that only a very small portion of signatures have class operators
and limited use of Kleene stars. This leads to a 
``bounded matching Regex" fragment (PM-Reg) we introduce later.
{\em To verify:
Using PM-Reg to approximate the general regex patterns
has no false positives/negatives on all the data set that we 
have experimented with.}

{\bf Efficient Application of Lookup argument:}
In \cite{Zkreg} it is observed that the generation of transition set
and check if they are valid can be separated, i.e.,
a valid encoding of a transition can be looked up in a pre-processed
transition table. We further
extend this approach and replace the zk-set component in \cite{Zkreg}
with a pre-processed lookup argument \cite{cq}, which leads to
prover complexity completely {\em independent of} 
the specification size, i.e., depending on the input string
length only. Compared with related work such as
\cite{Zombie23,ZkMiddleBox,Luo23} this approach largely improves
the practicality of $\zkregex$ in the context of non-membership proofs
with large automaton specification.

{\bf Scale-up to Multiple Patterns:} Given the observation that
a benign string has only a few ``near misses" to large collection
of regex patterns. We design a 2-stage zk-proof scheme for
disjunctive proof. For each regex signature, we collect 
a set of fixed patterns as its ``critical fixed patterns",
missing any of these patterns means that the input string
will not match the given pattern. We then compile an AC-DFA
over all such fixed patterns, and feed the input string
through the automaton. We collect a much small list of
final states that appear along the acceptance path, and
use these final states to retrieve the list of signatures
that should be proved (and the others be ignored).
For signatures of different types (e.g., PM-Reg and
general regular expressions), they are compiled into separate
circuits and folded as an IVC scheme by leveraging
ProtoGalaxy \cite{ProtoGalaxy}. The verifier has no knowledge
that whic patterns are actually being discharged, and via
which approach.

{\bf Bounded Matching Regex:} Based on data analysis
we present a {\em strictly weaker} fragment of regex,
however, which is capable of capturing $97\%$ of
ClamAV patterns. We develop efficient constraint
encoding techniques which encode the distance between
regex component and leave the fixed components to be 
handled by more efficient AC-DFA. All PM-Reg are compiled
into one general purpose circuit, while its
distance encoding constraints are loaded at run-time,
again leveraging lookup arguments.

{\bf Cryptographic Constructions:} Of independent interest,
we present a bounded zero knowledge version of the
$\cq$ lookup argument, which does not rely on the
Aurora lemma \cite{aurora19} and results in faster concrete
prover performance.  We present a modified
ProtoGalaxy folding scheme, enhanced with
folding techniques for zk-$\cq$ lookup argument.
We implement the schemes mentioned above fully, by
instrumenting a PSE-fork of the Halo2 system \cite{Halo2}.
We evaluate our scheme
over the full ClamAV signature set (which consists of
$38k$ regex signatures), and the data sets of emails and SQL 
database logs. 
The Halo2 system enhanced with $\cq$-lookups and ProtoGalaxy schemes
is general purpose, and can be applied to other problems.

